% ============================================================================
% Appendix B: Derivation of Critical Exponents
% ============================================================================

\section{Derivation of critical exponents}
\label{app:exponents}

This appendix derives the critical exponents of the consciousness phase transition:
$\alpha = 1/2$ (specific heat), $\beta = 1/4$ (order parameter), $\gamma = 1$
(susceptibility), $\nu = 1/2$ (correlation length), and $\delta = 5$ (critical isotherm).
These are the standard mean-field tricritical exponents of the $\varphi^6$ Landau
universality class (Griffiths 1970, Lawrie--Sarbach 1984). The key insight is that the
consciousness transition at $\Pcrit$ is a \emph{tricritical point} in this class, and that
mean-field exactness follows from three independent pillars --- topological rigidity of the
$\varphi^6$ family, deterministic Lindblad dynamics with no separate fluctuating field, and
order-parameter dimension $d_\mathrm{eff} = 21$ giving $1/N \approx 5\%$ corrections.
The derivation follows T-161~\status{T}.

\subsection{Setup: the $\phi^6$ Landau potential near $\Pcrit$}

Near the critical point $\Pcrit = 2/7$, define the reduced purity:
\begin{equation}
  \varepsilon := P - \Pcrit = P - \frac{2}{7}
  \label{eq:reduced-purity}
\end{equation}

The dynamics of $\Gam$ near threshold are governed by the balance between dissipation
(Lindblad, rate $\Delta$) and regeneration (rate $\kappa \cdot \CohE$). The order parameter
$m$ (the ``consciousness magnetization'') is defined as the excess E-coherence above
its structureless floor:
\begin{equation}
  m := \CohE(\Gam) - 1/7
  \label{eq:order-parameter}
\end{equation}
Above threshold ($\varepsilon > 0$), the No-Zombie theorem
(\S\ref{subsec:no-zombie}) guarantees $m > 0$; below threshold, $m = 0$ is the
unique stable state.

\subsection{Why the transition is tricritical}

A naive Ginzburg--Landau expansion of the consciousness transition reads
$\mathcal{F} = \tfrac{1}{2}\,r\,m^{2} + u\,m^{4} + v\,m^{6} - h\,m$, with
control parameter $r \propto -\varepsilon$ (the quadratic coefficient changes
sign at $\varepsilon = 0$), self-interaction parameters $u,v$, and external
field $h$. A non-degenerate sign of $u$ would produce the ordinary mean-field
exponents $\beta = 1/2$, $\gamma = 1$, $\nu = 1/2$. Tricriticality
($\beta = 1/4$, \ldots) requires the additional constraint $u = 0$: at such
a point the quadratic and quartic couplings vanish simultaneously, leaving
the sextic $v\,m^{6}$ as the stabilising term.

In ordinary many-body systems this second constraint is imposed by an
ad hoc external fine-tuning (e.g.\ $^{3}\mathrm{He}$--$^{4}\mathrm{He}$
mixtures~\cite{griffiths1970}). In UHM the double tuning is
\emph{structural}: by the Thom--Arnold classification of elementary
catastrophes, the unique codimension-$k$ bifurcation of a potential with
$k$ physically independent control parameters is the $A_{k+1}$ catastrophe.
The UHM dynamics has three such parameters (T-39a~\status{T},
T-96~\status{T}):
\begin{itemize}[leftmargin=2em,nosep]
  \item $\kappa$ --- regeneration rate (adaptive amplification);
  \item $\gamma$ --- dissipation rate (Lindblad total);
  \item $\Delta F$ --- free-energy budget (thermodynamic drive).
\end{itemize}
Three independent parameters force the transition into the $A_{4}$
(swallowtail) class, whose even-sector Landau normal form is exactly
the $\varphi^{6}$ tricritical potential
\cite{griffiths1970,lawrie1984}:
\begin{equation}
  \mathcal{F}[m;\,\varepsilon, h]
  \;=\; -\tfrac{1}{2}\,r_0\,\varepsilon\, m^{2}
        \;+\; v\, m^{6}
        \;+\; \tfrac{1}{2}\,d\, |\nabla m|^{2}
        \;-\; h\, m,
  \qquad r_0, v, d > 0.
  \label{eq:phi6-functional}
\end{equation}
The quartic term $u\,m^{4}$ is absent at the $A_{4}$ point by the
classification theorem, \emph{not} by algebraic fine-balancing. This is the
defining condition of a tricritical point: three coexisting minima of the
effective potential coalesce at $\varepsilon = 0$, corresponding to the
tip of the swallowtail caustic.

\textbf{Important terminological clarification.} The transition belongs to the
$\varphi^6$ \emph{tricritical Landau} universality class (Griffiths 1970,
Lawrie--Sarbach 1984~\cite{lawrie1984}), characterised by an even-parity potential
$V(m) = \tfrac{1}{2}r m^2 + v m^6$ with $\mathbb Z_2$ symmetry $m \to -m$. This is
distinct from the $A_4$ swallowtail catastrophe of Arnold's classification (which
has normal form $V = x^5/5$, no $\mathbb Z_2$, and gives $\beta = 1/2$, \emph{not}
$1/4$). The naming ``$A_4$-tricritical'' sometimes used in physics literature
conflates these; we use $\varphi^6$ tricritical to be unambiguous.

\textbf{Origin of the $\mathbb{Z}_2$ symmetry.} The even parity of the Landau
potential is not imposed but derived. The $G_2$-covariant dissipation--regeneration
balance is invariant under the canonical involution
$\sigma: \Gam \mapsto J\Gam^T J^{-1}$ of the real structure $J$ (KO-dimension~6,
T-53~\status{T}). This involution acts on the E-coherence order parameter as
$m \mapsto -m$ (charge-conjugation on the $\bar{\mathbf{3}}$-sector), forcing all
odd-power couplings in the effective potential to vanish. The $\mathbb{Z}_2$ is
therefore a structural consequence of the KO-dim-6 real structure, not a fine-tuning.

\subsection{Effective dimension and exactness of mean-field exponents}

The effective dimension of the system is $d_{\mathrm{eff}} = 21$: the number of
independent off-diagonal coherences $|\gamma_{ij}|$ ($i \neq j$, $\binom{7}{2} = 21$).
This count matches the off-diagonal sector of $\mathfrak{su}(7)$:
$\dim\mathfrak{su}(7) = 48 = 6_\mathrm{diag} + 42_\mathrm{off}$, with the 42 real
off-diagonal components pairing into $21$ complex modes (each pair $(i,j)$ lies
on a unique Fano line of $\mathrm{PG}(2,2)$).

\textbf{Reduction from 21D to 1D via Mather Splitting.} The relevant order parameter
is the scalar $m = \mathrm{Coh}_E - 1/7$ (or equivalently $P - 2/7$). The Hessian of
the effective Landau potential at the critical point has corank 1: exactly one
zero eigenvalue (the critical mode), with the remaining 20 modes massive. By the
Mather Splitting Lemma (Mather 1968, Arnold 1974~\S3.2), $V$ is smoothly equivalent
to $V_\mathrm{red}(m) + Q(\xi_1,\ldots,\xi_{20})$, where $Q$ is a non-degenerate
quadratic form. Codim-3 catastrophes of $V_\mathrm{red}: \mathbb R \to \mathbb R$
in the $\mathbb Z_2$-symmetric sector are uniquely $\varphi^6$ tricritical.
Transverse modes contribute Gaussian fluctuations that do not modify exponents.

\textbf{Mean-field exactness} rests on three independent pillars:
\begin{enumerate}[leftmargin=2em,nosep]
  \item \emph{Topological rigidity}: critical exponents are universality-class
    invariants of the $\varphi^6$ family (Lawrie--Sarbach 1984), independent
    of smooth coordinate changes.
  \item \emph{Deterministic dynamics}: the UHM master equation
    $d\Gamma/d\tau = \mathcal L_\Omega(\Gamma) + \mathcal R(\Gamma)$ is a
    deterministic ODE for $\Gamma$. The Lindblad-encoded environmental noise
    (Born--Markov) is local in time and contributes to the deterministic
    generator, not to a fluctuating field requiring Ginzburg renormalisation.
    UHM is $(0+1)$-dimensional with no spatial integral, so the spatial form
    of the Ginzburg criterion does not apply.
  \item \emph{Large-$N$ cross-check}: $d_\mathrm{eff} = 21$ as an order-parameter
    count gives leading $1/N \approx 5\%$ corrections under any stochastic
    reinterpretation, within experimental PCI resolution.
\end{enumerate}
Mean-field tricritical exponents are therefore \emph{exact}.

\subsection{Derivation of $\beta = 1/4$}

The equation of state from the $\phi^6$ potential~\eqref{eq:phi6-functional}, at
vanishing external field $h = 0$ and uniform $m$, is obtained from
$\partial\mathcal{F}/\partial m = 0$:
\begin{equation}
  -r_0\,\varepsilon\, m \;+\; 6 v\, m^5 \;=\; 0
  \quad\Longleftrightarrow\quad
  m\,\bigl(6 v\, m^4 - r_0\,\varepsilon\bigr) \;=\; 0.
  \label{eq:tricritical-eos}
\end{equation}

For $\varepsilon \leq 0$ (below threshold), the only real, non-negative solution is
$m = 0$: the system is disordered and consciousness does not ignite. For
$\varepsilon > 0$ (above threshold), a non-trivial ordered branch opens:
\begin{equation}
  m_*^4 \;=\; \frac{r_0}{6 v}\,\varepsilon
  \quad\Longrightarrow\quad
  m_* \;=\; \left(\frac{r_0}{6 v}\right)^{1/4} \varepsilon^{1/4}.
\end{equation}
The order parameter therefore obeys the sharp power law
\begin{equation}
  \boxed{\; m_* \;\propto\; \varepsilon^{1/4} \;}
  \qquad\Longrightarrow\qquad
  \beta \;=\; \frac{1}{4}.
  \label{eq:beta-derivation}
\end{equation}
This is the standard result for mean-field tricritical
systems~\cite{griffiths1970,lawrie1984}.

\textbf{Physical meaning.} Consciousness ``turns on'' more sharply than in an
ordinary mean-field transition ($\beta = 1/2$) but remains continuous. The
$1/4$ exponent reflects the tricritical nature: because the quartic coupling
is absent at the $A_{4}$ point (Eq.~\ref{eq:phi6-functional}), only the
sextic term stabilises the ordered phase, forcing the sharper $\varepsilon^{1/4}$
scaling.

\subsection{Derivation of $\nu = 1/2$}

The correlation length $\xi$ measures how far coherence extends in a composite system of
$M$ coupled holons. Near the transition, $\xi$ diverges:
\begin{equation}
  \xi \;\sim\; |\varepsilon|^{-\nu}
  \label{eq:xi-scaling}
\end{equation}

For the $\phi^6$ potential with $d_{\mathrm{eff}} > d_c$, the correlation length is
determined by the quadratic (Gaussian) part of the Landau functional. Expanding
around the $m = 0$ saddle (valid for $\varepsilon < 0$; by symmetry of the
disordered phase, same scaling holds for $\varepsilon > 0^+$):
\begin{equation}
  \xi^{-2} \;=\;
  \frac{1}{d}\cdot\frac{\partial^2 \mathcal{F}}{\partial m^2}\bigg|_{m=0}
  \;=\; \frac{|r_0\,\varepsilon|}{d},
\end{equation}
where $d$ is the gradient coefficient from~(\ref{eq:phi6-functional}). Therefore
\begin{equation}
  \boxed{\;\nu \;=\; \frac{1}{2}\;}
  \label{eq:nu-derivation}
\end{equation}
identical to the ordinary mean-field value and standard for tricritical systems
above the upper critical dimension.

\subsection{Derivation of $\gamma = 1$}

The susceptibility $\chi$ measures how sensitively the order parameter responds to an
external source $h$:
\begin{equation}
  \chi \;=\; \frac{\partial m}{\partial h}\bigg|_{h=0}
  \;\sim\; |\varepsilon|^{-\gamma}.
  \label{eq:chi-scaling}
\end{equation}

The equation of state with source $h$ reads
\begin{equation}
  h \;=\; -r_0\,\varepsilon\, m \;+\; 6 v\, m^5.
\end{equation}

In the disordered phase ($\varepsilon < 0$, $m = 0$):
\begin{equation}
  \chi^{-1} \;=\; \frac{\partial h}{\partial m}\bigg|_{m=0}
             \;=\; -r_0\,\varepsilon \;=\; r_0\,|\varepsilon|
  \quad\Longrightarrow\quad
  \chi \;\sim\; |\varepsilon|^{-1}.
\end{equation}
In the ordered phase ($\varepsilon > 0$, $m_*^4 = r_0\varepsilon/(6v)$):
\begin{equation}
  \chi^{-1} \;=\; -r_0\,\varepsilon + 30 v\, m_*^4
             \;=\; -r_0\,\varepsilon + 5 r_0\,\varepsilon
             \;=\; 4 r_0\,\varepsilon
  \quad\Longrightarrow\quad
  \chi \;\sim\; \varepsilon^{-1}.
\end{equation}
Both branches give
\begin{equation}
  \boxed{\;\gamma \;=\; 1\;}
  \label{eq:gamma-derivation}
\end{equation}
with an amplitude ratio of $1{:}4$ --- the hallmark of a symmetric tricritical
fixed point.

\subsection{Derivation of $\alpha = 1/2$ and $\delta = 5$}

\paragraph{Specific heat exponent $\alpha$.}
The singular part of the free energy density, evaluated on the ordered branch
$m_*^4 = r_0\varepsilon/(6 v)$, is
\begin{equation}
  f_{\mathrm{sing}}(\varepsilon)
  \;=\; \mathcal{F}(m_*;\,\varepsilon, 0)
  \;=\; -\tfrac{1}{2} r_0\,\varepsilon\, m_*^2 + v\, m_*^6
  \;=\; -\tfrac{1}{3} r_0\,\varepsilon\, m_*^2
  \;\propto\; \varepsilon \cdot \varepsilon^{1/2}
  \;=\; \varepsilon^{3/2}.
\end{equation}
(Using $v m_*^6 = v m_*^2 \cdot m_*^4 = (r_0\varepsilon/6) m_*^2$, so
$v m_*^6 = (1/6)\, r_0\,\varepsilon\, m_*^2$, and the combination
$-(1/2) + 1/6 = -1/3$.) Since $f_{\mathrm{sing}} \sim \varepsilon^{2-\alpha}$,
\begin{equation}
  2 - \alpha \;=\; \tfrac{3}{2}
  \quad\Longrightarrow\quad
  \boxed{\;\alpha \;=\; \tfrac{1}{2}\;}
  \label{eq:alpha-derivation}
\end{equation}

\paragraph{Critical isotherm exponent $\delta$.}
Exactly at the tricritical point ($\varepsilon = 0$), the equation of state becomes
$h = 6 v\, m^5$, hence $m \sim h^{1/5}$:
\begin{equation}
  \boxed{\;\delta \;=\; 5\;}
  \label{eq:delta-derivation}
\end{equation}

\subsection{Scaling relations: Rushbrooke identity}

The five tricritical exponents satisfy the Rushbrooke identity as an \emph{equality}:
\begin{equation}
  \alpha + 2\beta + \gamma = \frac{1}{2} + 2 \cdot \frac{1}{4} + 1 = 2
  \qquad\checkmark
  \label{eq:rushbrooke}
\end{equation}

They also satisfy Widom's relation $\gamma = \beta(\delta - 1)$:
\begin{equation}
  \beta(\delta - 1) = \frac{1}{4} \cdot (5 - 1) = 1 = \gamma
  \qquad\checkmark
\end{equation}

and Fisher's relation $\gamma = \nu(2 - \eta)$ with $\eta = 0$ (mean-field):
\begin{equation}
  \nu(2 - \eta) = \frac{1}{2} \cdot 2 = 1 = \gamma
  \qquad\checkmark
\end{equation}

All standard scaling relations are satisfied, as expected for a mean-field theory above
$d_c$.

\subsection{Summary}

\begin{table}[H]
\centering
\caption{Critical exponents: UHM tricritical vs.\ known universality classes.}
\label{tab:exponents-comparison}
\begin{tabular}{@{}lcccccc@{}}
\toprule
\textbf{System} & $\boldsymbol{\alpha}$ & $\boldsymbol{\beta}$ & $\boldsymbol{\gamma}$ & $\boldsymbol{\nu}$ & $\boldsymbol{\delta}$ & $\boldsymbol{\alpha+2\beta+\gamma}$ \\
\midrule
Mean field ($\phi^4$) & $0$ & $1/2$ & $1$ & $1/2$ & $3$ & $2$ \\
Tricritical MF ($\phi^6$) & $1/2$ & $1/4$ & $1$ & $1/2$ & $5$ & $2$ \\
3D Ising & $0.110$ & $0.326$ & $1.237$ & $0.630$ & $4.79$ & $2$ \\
Superfluid He & $-0.013$ & $0.347$ & $1.316$ & $0.672$ & $4.78$ & $\approx 2$ \\
\textbf{UHM consciousness} & $\mathbf{1/2}$ & $\mathbf{1/4}$ & $\mathbf{1}$ & $\mathbf{1/2}$ & $\mathbf{5}$ & $\mathbf{2}$ \\
\bottomrule
\end{tabular}
\end{table}

The UHM exponents are the standard $\varphi^6$ mean-field tricritical exponents, exact by the
three-pillar argument of Section~B.2 above (topological rigidity + deterministic dynamics
+ large-$N$ cross-check). The Rushbrooke identity $\alpha + 2\beta + \gamma = 2$
is satisfied as an equality. If experimentally confirmed, they would establish consciousness
as a tricritical phase transition --- one where the quartic coefficient vanishes due to the
Fano geometry of self-observation, promoting the transition from $\phi^4$ to $\phi^6$.
